\section{Extended Feasibility Rubric (Detailed)}
\label{sec:feasibility_rubric_section}

\subsection{Smart-Contract Feasibility Rubric}
\label{sec:feasibility_rubric_full}

\subsubsection{Why feasibility must be a first-class criterion}
This thesis evaluates multiple candidate metrics for evidence-based developer reputation.
However, the project explicitly requires that feasibility as a smart contract is considered as a
selection criterion and that a proof-of-concept (PoC) is implemented for one chosen metric.
The project scope further constrains the design through a conceptual evidence workflow where developers submit
a signed form containing a link and a hash of an archive, the artifact and metric are verified, compact evidence is written
to the ledger, and a sibling search/index system can observe each new data point. Because Fabric chaincode cannot retrieve an
external archive or traverse Git, the implemented interpretation assigns retrieval and metric computation to a deterministic
off-chain verifier and reserves bounded authorization, validation, commitment, state, and event operations for chaincode.
Therefore, feasibility here is not a generic ``blockchain compatibility'' notion, but the practical
ability to implement (or safely approximate) the required pipeline under smart-contract constraints
while preserving auditability and integrity.

In the blockchain literature, feasibility is shaped by system constraints that include (i) performance
limits and throughput constraints, (ii) storage growth and distribution overhead, and (iii) challenges
in connecting on-chain logic to external services and data (oracle-style issues). These constraints
make heavy computation and large state growth poor candidates for on-chain implementation, often
motivating hybrid architectures in which bulk data remains off-chain and the chain stores commitments
and compact metadata~\cite{blockchainsurvey}. Smart-contract ecosystems add further constraints: execution is costed (gas),
code and execution patterns are shaped by platform limits, and determinism is required.
Empirical work on Solidity metrics shows that many metric distributions are truncated and exist in
restricted ranges due to platform constraints, which supports treating feasibility as a distinct and
non-trivial selection criterion rather than assuming traditional software metrics transfer
unchanged~\cite{smartmetrics2023}.

\subsubsection{Definition: smart-contract feasibility}
This thesis defines the \textbf{smart-contract feasibility} of a candidate metric as:
\begin{quote}
The ability to compute, verify, and store a metric (or its verified output) within the defined evidence workflow
with bounded on-chain computation and storage, minimal reliance on unverifiable external inputs, strong auditability
(re-checking from archived evidence), and explicit binding to developer identity/provenance.
\end{quote}

This definition is intentionally \emph{architecture-aware}. A candidate can be feasible in one of two ways:
\begin{itemize}[leftmargin=*]
  \item \textbf{On-chain computable feasibility:} the metric can be computed directly in the contract with small, bounded logic.
  \item \textbf{Hybrid feasibility (recommended for the PoC):} the metric is computed off-chain from an artifact whose hash and link are verified,
  and the contract stores only compact metric values plus hash commitments to evidence/inputs to enable auditability.
\end{itemize}
The rubric below evaluates feasibility in a way that allows both patterns, while explicitly penalizing designs that
require unverifiable oracles, frequent high-throughput writes, or large on-chain state.

\subsubsection{Rubric structure and scoring}
Each metric receives a score on 10 dimensions, each scored in $\{0,1,2\}$ (0 = poor feasibility, 2 = strong feasibility).
The rubric is used \textbf{only after} the multi-metric comparative analysis, to select the single PoC metric.

\paragraph{Rubric application procedure.}
\begin{enumerate}[leftmargin=*]
  \item For each candidate metric $M_i$, record a one-paragraph implementation account specifying the inputs,
  computation, stored values, update frequency, and attesting party.
  \item Score each rubric dimension using the guidance below.
  \item Compute a total feasibility score and also record \emph{exclusion conditions} (dimensions where score = 0
  implies unacceptable risk for the PoC).
  \item Select the PoC metric as the best trade-off between (i) comparative performance under C1--C4 and
  (ii) feasibility score, with tie-breaking favoring determinism, auditability, and bounded cost.
\end{enumerate}

\subsubsection{Feasibility dimensions (detailed)}
\label{sec:feasibility_dimensions_detailed}

\begin{table}[H]
\centering
\scriptsize
\renewcommand{\arraystretch}{1.14}
\begin{tabularx}{\textwidth}{p{2.8cm}p{0.8cm}X}
\toprule
\textbf{Dimension} & \textbf{0--2} & \textbf{Scoring rule (detailed)} \\
\midrule
D1. Determinism of computation & 0--2 &
\textbf{2:} metric computation is deterministic given the archived artifact and explicit parameters (window definition, tool version).
\textbf{1:} deterministic only under additional assumptions (toolchain pinned; minor nondeterminism mitigated).
\textbf{0:} depends on nondeterministic factors or ambiguous extraction (e.g., floating-point randomness, unstable parsing). \\

D2. Bounded on-chain computation & 0--2 &
\textbf{2:} on-chain logic is constant-time or small bounded checks (signature format check, hash equality, record writes).
\textbf{1:} moderate bounded computation per update (limited loops over small arrays).
\textbf{0:} requires heavy analysis on-chain (AST parsing, full-history scanning) or unbounded loops. \\

D3. Bounded on-chain storage growth & 0--2 &
\textbf{2:} store only compact values + commitments (hashes) + minimal metadata.
\textbf{1:} store moderate aggregates (rolling sums, bounded windows).
\textbf{0:} requires storing large histories/artifacts or per-commit detailed logs on-chain. \\

D4. Update frequency feasibility & 0--2 &
\textbf{2:} metric can be updated per release or per time window (low write rate).
\textbf{1:} requires medium-frequency updates (e.g., per PR/issue) but can be batched.
\textbf{0:} requires high-frequency updates (per commit/event stream) and becomes throughput-limited. \\

D5. External data/oracle dependence & 0--2 &
\textbf{2:} inputs are self-contained in the archived artifact + signed form, verifiable by hash commitments.
\textbf{1:} uses off-chain data but can be posted as signed attestations with clear trust model.
\textbf{0:} requires unverifiable external services or subjective judgments (strong oracle problem). \\
\bottomrule
\end{tabularx}
\caption{Feasibility dimensions with detailed scoring rules (D1--D5).}
\label{tab:feasibility_dimensions}
\end{table}

\begin{table}[H]
\centering
\scriptsize
\renewcommand{\arraystretch}{1.14}
\begin{tabularx}{\textwidth}{p{2.8cm}p{0.8cm}X}
\toprule
\textbf{Dimension} & \textbf{0--2} & \textbf{Scoring rule (detailed)} \\
\midrule
D6. Auditability and re-checking & 0--2 &
\textbf{2:} third parties can recompute metric from archive retrieved via link and verified by on-chain hash; measurement report is committed.
\textbf{1:} reproducible with effort (tooling dependencies; partial evidence).
\textbf{0:} cannot be reliably reproduced or verified from public evidence. \\

D7. Identity binding and provenance & 0--2 &
\textbf{2:} metric evidence is bound to developer identity via signed submissions (and the same identity used consistently across submissions).
\textbf{1:} partial binding (mapping from commit authors; weaker identity guarantees).
\textbf{0:} attribution cannot be trusted or can be trivially spoofed. \\

D8. Integrity of artifact linkage & 0--2 &
\textbf{2:} strong commitment: archive hash stored on-chain and verified against retrieved archive; versioning explicit.
\textbf{1:} link is stable but hash commitment incomplete or relies on weak assumptions.
\textbf{0:} artifact cannot be pinned; link rot or mutable artifact breaks verification. \\

D9. Metric-value representation (range + precision) & 0--2 &
\textbf{2:} metric stored as small integers or fixed-point with clear bounds; low risk of overflow/precision issues.
\textbf{1:} requires careful normalization/scale but feasible (e.g., store numerator/denominator separately).
\textbf{0:} requires floating point or complex statistics not representable safely on-chain. \\

D10. Adversarial resilience in the configured ledger context & 0--2 &
\textbf{2:} metric is hard/expensive to inflate without leaving detectable traces, especially with hash-anchored evidence.
\textbf{1:} partially gameable but mitigations exist (normalization, batching, cross-check with other metrics in study).
\textbf{0:} trivially gameable with low-cost manipulations (e.g., verbosity inflation) and weak detectability. \\
\bottomrule
\end{tabularx}
\caption{Feasibility dimensions with detailed scoring rules (D6--D10).}
\label{tab:feasibility_dimensions_cont}
\end{table}

\subsubsection{Evidence checklist for each dimension}
To keep feasibility rigorous and evidence-based, each dimension must be supported by explicit evidence:
\begin{itemize}[leftmargin=*]
  \item \textbf{D1 Determinism:} exact inputs (archive hash, link, and version identifier), the window definition, and, for off-chain computation,
  the pinned tool version and configuration; an explanation of how the PoC ensures that the same artifact yields the same output.
  \item \textbf{D2/D3 Cost bounds:} the contract operations required per data point (comparisons and storage writes), together with a boundedness justification.
  \item \textbf{D4 Update frequency:} the update cadence (per release or time window) and a justification for avoiding per-commit updates.
  \item \textbf{D5 Oracle dependence:} any off-chain sources beyond the archive and, where applicable, whether they are hash-anchored or signed attestations.
  \item \textbf{D6 Auditability:} the \emph{re-check protocol}: archive retrieval, hash verification, metric recomputation, and comparison with the ledger value.
  \item \textbf{D7/D8 Provenance and linkage:} an explanation of how the signed form and hash commitment prevent artifact substitution.
  \item \textbf{D9 Representation:} the encoding (integer, fixed-point, or separate numerator and denominator), overflow handling, and normalization.
  \item \textbf{D10 Adversarial model:} low-cost manipulations and whether evidence-consistency checks can detect them.
\end{itemize}

\subsubsection{Feasibility exclusion conditions}
Even if a metric is strong in the comparison, it is a weak PoC candidate if it triggers exclusion conditions:
\begin{itemize}[leftmargin=*]
  \item \textbf{Exclusion Condition A:} score 0 on D1 (non-deterministic), D6 (non-auditable), or D8 (artifact cannot be pinned).
  \item \textbf{Exclusion Condition B:} requires high-frequency updates (D4=0) that exceed realistic ledger throughput.
  \item \textbf{Exclusion Condition C:} depends on unverifiable external facts (D5=0) without a signed-attestation model.
\end{itemize}

\subsubsection{Mapping to the project evidence workflow}
The conceptual project workflow comprises signature checking, archive retrieval, hash verification, metric measurement,
ledger recording, and indexing notification. The rubric operationalizes this in dimensions D1, D6, D7, and D8
(determinism, auditability, identity binding, artifact linkage), and constrains practical implementation via D2--D5 and D9
(bounded compute/storage, update frequency, oracle dependence, representation). Section~\ref{sec:phase3_hybrid_boundary}
records the implemented split: archive and Git operations are off-chain, whereas chaincode performs bounded validation,
commitment, state, and event operations.

\subsubsection{Application to the five thesis candidates}
\label{sec:rubric_template_apply}
Table~\ref{tab:rubric_template} records the scoring shape used for each candidate. The actual scored application is reported
after C1--C4 in Section~\ref{sec:phase2_feasibility_scoring}, where each dimension receives a value and rationale.

\begin{table}[htbp]
\centering
\footnotesize
\renewcommand{\arraystretch}{1.12}
\resizebox{\textwidth}{!}{%
\begin{tabular}{lcccccccccc}
\toprule
Metric & D1 & D2 & D3 & D4 & D5 & D6 & D7 & D8 & D9 & D10 \\
\midrule
M1 Snapshot size (LOC/SLOC) & N/A & N/A & N/A & N/A & N/A & N/A & N/A & N/A & N/A & N/A \\
M2 Patch complexity (churn) & N/A & N/A & N/A & N/A & N/A & N/A & N/A & N/A & N/A & N/A \\
M3 Activity cadence & N/A & N/A & N/A & N/A & N/A & N/A & N/A & N/A & N/A & N/A \\
M4 Ownership/concentration & N/A & N/A & N/A & N/A & N/A & N/A & N/A & N/A & N/A & N/A \\
M5 Maintenance responsiveness (design ext.) & N/A & N/A & N/A & N/A & N/A & N/A & N/A & N/A & N/A & N/A \\
\bottomrule
\end{tabular}
}
\caption{Feasibility rubric scoring shape. Actual scores are reported in Section~\ref{sec:phase2_feasibility_scoring}.}
\label{tab:rubric_template}
\end{table}





